\chapter{猜数字}
\label{ch:02}

用一个交互循环串起输入、比较和分支。答案先固定为 42，使输入与输出能够重复核对。随机数属于生成测试数据的方法，不是理解循环的前提。

\section{把交互拆成步骤}

每轮先读取一个整数，再与答案比较：较小输出 small，较大输出 large，相等输出 found 并结束循环。答案在循环外声明，猜测值在每轮循环内创建。这个位置差异决定了变量可以在哪些语句中使用。

\section{实现并运行}

下面的完整程序使用 inputInt 读取整数。依次输入 30、60、42；每个数字后按回车。下方输出不包含终端对输入字符的回显。

\par
\medskip
\begin{lstlisting}[style=chali]
void main()
{
    int64 answer = 42;
    while (true)
    {
        int64 guess = inputInt();
        if (guess < answer)
        {
            println("small");
        }
        elif (guess > answer)
        {
            println("large");
        }
        else
        {
            println("found");
            break;
        }
    }
}
\end{lstlisting}
\par
\medskip

\par
\medskip
\begin{lstlisting}[style=out]
small
large
found
\end{lstlisting}
\par
\medskip

\section{分支与状态变化}

if、elif、else 形成一条分支链，一轮选择其中一个分支。else if 是 elif 的语法糖。break 退出当前循环，随后执行循环后面的语句。若需要统计次数，在循环外声明计数器，并在读取成功后的固定位置更新它。

while(true) 的结束依赖显式退出路径。移除 break 后，程序会继续请求输入；测试时应提供能够覆盖退出路径的数据。

\section{把输入错误留到合适的层处理}

这里假定输入符合整数读取接口要求。面向用户的文本工具通常先读取整行，再用 std::parse 解析，并通过 Result 处理格式错误与输入结束。第 12 章讨论文本，第 16 章讨论错误返回。不要把“输入无效”和“猜测不正确”合并成同一种状态。随机答案可查\href{https://github.com/rabbitc999/ChaLi/blob/master/doc/library/README.md}{随机数接口}，测试时仍保留固定答案用例。

\section{习题}

\begin{enumerate}

\item 输入 42，输出哪些行？

\item 增加猜测次数。

\item 用 else if 改写分支。

\item 为什么答案放在循环外？

\end{enumerate}

\subsection*{参考解答与验收要点}

\begin{enumerate}

\item 只输出 found。

\item 循环外初始化，在每轮读取后加一，相等分支输出次数。

\item 分支条件和顺序保持一致，结果不变。

\item 答案需要跨迭代保存；每轮猜测只在本轮使用。

\end{enumerate}

下一章：\hyperref[ch:03]{第 3 章}。

\section{番外：龟兔猜数字}

改写自伊索寓言《龟兔赛跑》。

兔子常常讥笑乌龟猜得慢。乌龟不服，说：“你跑得虽快，也未必猜得准。目标一到一百，我们比一比。”狐狸做裁判，先把答案藏进一个整数里，又特别声明：输入字符串不算猜。

兔子跳到终端前，先猜九十九，再猜二，接着输入一个“很大的数”。解析器报错。兔子觉得自己已经进行了三轮高强度推理，便躺到树下，说剩下的七十多次机会足够睡个午觉。

乌龟先猜五十。裁判说小，它便删去一半候选；猜七十五，大，又删一半。每走一步，都核对边界，连同一个错误也不肯重犯。狐狸看得直打呵欠：这个算法连戏剧冲突都对半砍。

兔子醒来，见乌龟还在纸上写字，笑道：“这么久才猜了六次！”说罢又输入九十九。裁判答：“这个分支你早来过。”

乌龟此时已写下答案。兔子奔到终点，发现验收表只有一行：“正确，七次以内。”它气得跳脚，说自己的输入速度遥遥领先。狐狸点头，把“每秒错误次数第一名”的奖牌挂在了它脖子上。

\subsection*{技术对照}

输入失败与猜错是不同分支。先检查输入转换结果，再比较大小；上下界更新必须保留尚有可能的候选，并保证每轮取得进展。
